\section{Introduction}

Search-and-Rescue (SAR) missions represent a critical operational domain for unmanned aerial systems~\cite{silvagni2017multipurpose}. When survivors are stranded in vast, inaccessible terrain, emergency response windows are severely constrained by environmental conditions. Unmanned Aerial Vehicles (UAVs) provide rapid deployment, flexible sensor mobility, and elevated coverage.

MANAR (Multi-sensor Autonomous Navigation and Aerial Rescue) is a supervised-autonomy SAR system designed to balance human oversight against automated execution~\cite{goodrich2008towards}. The term \emph{Autonomous} refers specifically to deterministic navigation capabilities under supervisory human authority. Fully manual teleoperation imposes heavy cognitive workloads on operators, while unconstrained autonomous systems introduce failure risks during sensor dropouts or communication loss. In MANAR, human operators retain supervisory authority over mission configuration and target verification, while a C++ control core executes deterministic navigation, route optimization, and safety monitoring.

\subsection{Three-Tier Maturity Framework}
To maintain engineering rigor, MANAR's capabilities are categorized across three distinct maturity tiers, summarized in Table~\ref{tab:implementation_status}:
\begin{enumerate}
    \item \textbf{Implemented Functional Prototype:} Fully compiled C++ control logic, double-precision Haversine navigation, $O(n^2)$ greedy route optimization with operator override, lawnmower sector sweeps, discrete battery simulation, structured event logging, and an ONNX Runtime C++ visual human detector module (YOLO11n) benchmarked against a D-FINE-N reference baseline.
    \item \textbf{Selected Closed Architecture:} Authoritative state ownership in C++, React/TypeScript web interface over WebSocket + JSON, and on-demand Return-to-Launch (RTL) navigation semantics.
    \item \textbf{Planned Future Work:} Multi-frame candidate tracking (Kalman / ByteTrack), temporal multi-sensor evidence fusion (Mamba SSM), PX4 SITL/HITL flight integration, and field deployment.
\end{enumerate}

\begin{table}[t]
\caption{MANAR Capability Implementation Status}
\label{tab:implementation_status}
\centering
\footnotesize
\begin{tabular}{@{}lc@{}}
\toprule
\textbf{Capability} & \textbf{Status} \\
\midrule
Deterministic C++ control core & \ding{51} Implemented \\
Haversine navigation geometry & \ding{51} Implemented \\
Multi-location sector sequencing & \ding{51} Implemented \\
Greedy route optimization ($O(n^2)$) & \ding{51} Implemented \\
Operator optimization toggle ($Y/N$) & \ding{51} Implemented \\
Lawnmower search grid logic & \ding{51} Implemented \\
Payload component state controls & \ding{51} Implemented \\
Structured event logging & \ding{51} Implemented \\
Visual candidate detector (YOLO11n + ORT) & \ding{51} Implemented \\
Reference detector baseline (D-FINE-N) & \ding{51} Implemented \\
Empirical 26-video SAR benchmark & \ding{51} Implemented \\
\midrule
React / TypeScript Web GUI & Selected \\
WebSocket + JSON IPC transport & Selected \\
Authoritative C++ state ownership & Selected \\
\midrule
Candidate tracking \& Kalman filtering & Planned \\
Multimodal evidence fusion (Mamba SSM) & Planned \\
PX4 SITL / HITL simulation & Planned \\
Physical flight deployment & Planned \\
\bottomrule
\end{tabular}
\end{table}

\subsection{Key Contributions}
The main contributions of this work are:
\begin{itemize}
    \item Formulation and verification of an $O(n^2)$ greedy nearest-neighbor route optimizer using Haversine geodetics~\cite{sinnott1984haversine} with operator override privilege.
    \item Integration and empirical validation of an ONNX Runtime C++ visual human detector module (YOLO11n) achieving $29.2\text{ ms}$ latency with zero false positives on clutter across 26 SAR video datasets.
    \item Specification of a dual-regime adaptive telemetry rate policy and analytical radio power consumption model across transmission frequencies.
    \item Formalization of parameterized motion, two-level RF attention, and simulated battery failsafe progression models.
    \item Functional verification of search plan locking, sequential sector progression, and Return-to-Launch navigation preemption.
\end{itemize}
